% Generated draft: requires mathematical and source review.
\documentclass[11pt]{article}
\usepackage[margin=25mm]{geometry}
\usepackage{fontspec}
\setmainfont{Latin Modern Roman}
\usepackage{amsmath,amssymb,mathtools,unicode-math}
\setmathfont{Latin Modern Math}
\usepackage{enumitem,needspace,xurl,hyperref}
\setlength{\parindent}{0pt}
\setlength{\parskip}{5pt}
\setlength{\emergencystretch}{3em}
\newcommand{\N}{\mathbb{N}}
\newcommand{\Z}{\mathbb{Z}}
\newcommand{\Q}{\mathbb{Q}}
\newcommand{\R}{\mathbb{R}}
\newcommand{\C}{\mathbb{C}}
\newcommand{\F}{\mathbb{F}}
\newcommand{\A}{\mathbb{A}}
\newcommand{\PP}{\mathbb{P}}
\newcommand{\E}{\mathbb{E}}
\newcommand{\eps}{\varepsilon}
\newcommand{\dd}{\,\mathrm{d}}
\begin{document}
\section*{UnsolvedMath ID: 30006259}
\textbf{Source problem code:} catalogue entry
\subsection{Definitions and mathematical statement}
The supplied record does not specify the hypotheses on the field \(k\), its valuation ring \(\mathfrak{o}\), the group \(G\), or the parameter \(f\), nor does it give the construction of the object called the wonderful group scheme \(G_{x,f}\). The following question is therefore conditional on recovering those data and hypotheses from the cited source.

Write \(\mathcal{B}(G)\) for the Bruhat--Tits building intended in the source. Let \(C\subseteq\mathcal{B}(G)\) be a chamber containing \(x\), and put
\[
I:=G(k)_{C,0},\qquad X:=G_{x,f}(\mathfrak{o}),
\]
where the source identifies \(I\) as the associated Iwahori subgroup and \(X\) denotes the set of \(\mathfrak{o}\)-valued points. The convention for chambers and whether boundary points are permitted are not specified.

\emph{Orbit-classification question.} For the intended action of \(I\times I\) on \(X\), classify all orbits. More explicitly, after the action
\[
\alpha:(I\times I)\times X\longrightarrow X
\]
has been specified, determine a parameter set \(A\) and a complete invariant \(q:X\to A\) such that \(q\) is surjective and, for every \(u,v\in X\),
\[
q(u)=q(v)
\quad\Longleftrightarrow\quad
\exists(i_1,i_2)\in I\times I:\ v=\alpha((i_1,i_2),u).
\]
This is the requested analogue of the classification of \((B\times_k B)\)-orbits in the wonderful compactification mentioned in the excerpt. Neither \(B\) nor the action on \(X\) is defined in the supplied material; in particular, a left--right multiplication formula is not imposed here.

\emph{Stable-piece question.} Describe the diagonal-stable pieces of \(X\) intended by the source, in analogy with the stable-piece decomposition of the wonderful compactification. For a group \(H\), its diagonal subgroup is
\[
\operatorname{Diag}(H):=\{(h,h):h\in H\}\subseteq H\times H.
\]
The original excerpt specifically refers to \(G_{x,f}(\mathfrak{o})^{\operatorname{Diag}}\)-stable pieces in \(G_{x,f}(\mathfrak{o})\). It does not define the relevant diagonal action or the pieces. Thus no precise partition, parametrization, or equivalence with diagonal-subgroup orbits can be asserted from this record alone.

The excerpt motivates both questions by similarities with the wonderful compactification when \(f=0\). It describes its Theorems 1 and 3 as possible first steps toward a Bruhat--Tits theory for wonderful compactifications, and expresses an expectation of extending the results to more general wonderful varieties in the sense of its reference [4]. These are source qualifications and motivations, not additional established conclusions or specified extension statements.
\subsection{Short English statement}
For a chamber \(C\) containing \(x\), classify the orbits of the associated Iwahori subgroup product \(G(k)_{C,0}\times G(k)_{C,0}\) on \(G_{x,f}(\mathfrak{o})\), and describe the intended diagonal-stable pieces, by analogy with the wonderful compactification. The construction, hypotheses, actions, and meaning of stable pieces require clarification from the source before this becomes a fully specified mathematical problem.
\subsection{Sources}
\begin{itemize}
\item Local UnsolvedMath registry: supplied complete source record, permanent ID \(30006259\), title ``Orbit Decompositions in Wonderful Group Schemes,'' exact source problem code \texttt{catalogue entry}. The record belongs to the Oberwolfach Reports Open Problems collection, gives proposed year 2025, and names no individual proposer.
\item Supplied report citation: \emph{Algebraic Groups} (2025), Oberwolfach Reports, pp.~941--1010, DOI \texttt{10.4171/owr/2025/20}. \url{https://doi.org/10.4171/owr/2025/20}.
\item Supplied literature source, ``OWR 2025 report'': \url{https://ems.press/content/serial-article-files/51849}.
\item Supplied literature source, ``Li, wonderful embeddings for group schemes'': \url{https://arxiv.org/abs/2505.12777}. No fuller author attribution or bibliographic details are supplied.
\item Literature assessment dated \texttt{2026-08-22}: the supplied curation labels the problem open and states that the wonderful group-scheme construction is available while the requested Iwahori orbit classification and diagonal-stable-piece description remain open questions. This is supplied dated curation, not an independently verified literature assessment.
\item The registry labels the statement \texttt{corrected\_verified} and reports that the official report was checked and notation repaired. This is supplied verification metadata, not verification performed for this draft. The excerpt refers to [2, \S~2.1], [5], [4], and Theorems 1 and 3, but their bibliographic identities or statements are not included in the supplied record.
\end{itemize}
\end{document}
